\documentclass[10pt,a4paper]{amsart}
\usepackage[margin=19mm]{geometry}
\usepackage{amsmath,amssymb,mathtools}
\usepackage[scr=rsfs,cal=euler]{mathalfa}
\usepackage{xfrac}
\usepackage[unicode,hidelinks,bookmarks=false]{hyperref}
\newcommand{\ie}{i.e.\ }
\newcommand{\cf}{cf.\ }
\newcommand{\resp}{respectively\ }
\newcommand{\Subsection}{\subsection}
\providecommand{\nobreakdash}{\nobreak\hbox{-}}
\setlength{\emergencystretch}{2em}
\multlinegap0pt
\usepackage{amssymb}
\usepackage{mathtools}
\usepackage{tikz}
\newtheorem{theorem}{Theorem}
\newcommand{\N}{\mbox{${\mathbb N}$}}
\newcommand{\Z}{\mbox{${\mathbb Z}$}}
\title{Rainbow Trapezoids with Given Area}
\author{Sukumar Das Adhikari, T\'assio Naia, Oriol Serra}
\date{Source: arXiv:2603.13841v1 (CC BY 4.0)}
\begin{document}
\maketitle

\noindent\emph{Source and licence.} Rainbow Trapezoids with Given Area. Version arXiv:2603.13841v1 (CC BY 4.0), distributed by arXiv under the Creative Commons Attribution 4.0 International licence, http://creativecommons.org/licenses/by/4.0/, as recorded in the arXiv licence metadata for this exact version and shown on the abstract page https://arxiv.org/abs/2603.13841v1. Full licence notice: \texttt{licenses/arxiv-2603.13841-CC-BY-4.0.txt}.\par\smallskip

\noindent\textit{Editorial scope.} Counted unit: the article's theorem t:rb-triangle-vs-even-mon-rects (source0.tex lines 198--202). The article's complete original proof of it is delivered with it (source0.tex lines 204--269, one \texttt{proof} environment of 66 lines). No mathematics is written, repaired or re-derived: the statement and the proof are the article's own lines, in the article's own notation, and no retained line is substituted or re-flowed.  No further source-local context is needed: every cross-reference of every retained line resolves inside the two delivered spans.  Nothing was omitted from any retained span, so the package carries no omission record.  No retained line contains a citation, so the wrapper carries no bibliography and no bibliography entry.  No retained line contains a figure environment, an image-inclusion command or drawing code, so no image is delivered and the package declares no asset dependency.  Editorial formatting only, in addition: an AMS \texttt{amsart} preamble that restates the article's own declarations and macros used by the retained lines, namely the environments \texttt{theorem} on separate counters, together with the standard packages those lines need.  The retained environments print in this document as Theorem 1 in order of appearance, whereas the article prints the same items with the numbering of its own full document.  The complete native source of the article as distributed in the arXiv e-print for this exact version is bundled unchanged as \texttt{sources/196353/source0.tex}.
 \begin{theorem}\label{t:rb-triangle-vs-even-mon-rects}
      In every $r$-coloring of~$\Z^2$ with $r\ge 3$ colors there is either
      a rainbow triangle of area~$1/2$ or monochromatic rectangles
      of each even area whose sides are parallel to the axes.
    \end{theorem}

    \begin{proof}
      Fix an $r$-coloring $\chi$ of~$\Z^2\!$.
      In what follows (unless explicitly stated otherwise) by a \emph{ line} we mean a horizontal line;
      by a \emph{rainbow triangle} we mean the vertices of a triangle of area~$1/2$
      whose colors are pairwise distinct;
      and by a \emph{monochromatic rectangle} we
      mean a set of points of the same color which form the vertices of a rectangle
      whose sides are parallel to the axes of~$\Z^2\!$.

      \medskip
      \emph{Case 1.} Some line $\ell$ contains points of $3$~distinct colors
      (i.e., $\bigl|\chi(\ell)\bigr|\ge 3$).
      By renaming the colors if necessary, we may assume that there is a  pair  of consecutive points
      in $\ell$ colored with~$\{a,b\}$ and a second pair of consecutive points colored with  $\{a',b'\,\}$
with $\{a,b\}\neq \{a',b'\,\}$.
      Then any occurrence of a color not in $\{a,b\}$ or not in $\{a',b'\,\}$  in a line adjacent to~$\ell$
      produces a rainbow triangle of area $1/2$ with the first or the second pair respectively.
      We conclude that if a rainbow triangle is not present, then both lines adjacent to $\ell$ receive uniquely a color in $\{a,b\}\cap \{c,d\}$. Since this intersection has size at most~$1$, all points in the lines adjacent to~$\ell$ are assigned the same color, yielding monochromatic rectangles of every even area.

      \medskip
      \emph{Case 2.} Some pair of consecutive lines $\ell_1,\, \ell_2$ contain points of at least $3$ distinct
colors
      (i.e., $\bigl|\chi(\ell_1\cup\ell_2)\bigr|\ge 3$ for some pair of consecutive lines~$\ell_1,\,\ell_2$).
      We may assume none of the lines contains 3~colors, for otherwise we are done by Case~1.
      Without loss of generality, we may also assume that $\ell_1$ contains adjacent points of colors~$a,\,b$,
      and $\ell_2$ contains a point of color~$c$. These three points form a rainbow triangle.

      \medskip
      \emph{Case 3.} There exist three consecutive lines $\ell_1,\,\ell_2$, and~$\ell_3$ that
      together contain at least~$3$ colors
      (i.e., $\bigl|\chi(\ell_1\cup\ell_2\cup\ell_3)\bigr|\ge 3$).
      We may assume that $\{a,\,b,\,c\}\subseteq \chi(\ell_1\cup\ell_2\cup\ell_3)$, that neither
      of the previous cases applies (in particular, $\bigl|\chi(\ell_i)\bigr|\le 2$ for each $i\in\{1,\,2,\,3\}$),
      and that $\ell_2$ lies between $\ell_1$ and~$\ell_3$.

      If each of these lines contains a single color, then we obtain a rainbow triangle.\footnote{{In fact,
      rainbow triangles of every  area multiple of $1/2$}.}
      Otherwise, one of these lines contains two colors.
      We may assume that $\bigl|\chi(\ell_2)\bigr|=1$, %: indeed, if $\bigl|\chi(\ell_2)\bigr|>2$ we are done by Case~1,
      as if $\bigl|\chi(\ell_2)\bigr|=2$ then either $\ell_1$ or~$\ell_3$ contains
      a third color and we are done by Case~2.
      In particular, we may suppose without loss of generality that $\chi(\ell_1)=\{a,b\}$, that $\chi(\ell_2)=\{a\}$ and $c\in\chi(\ell_3)$.

      Note that if $\chi(\ell_3)\not\subseteq \{a,c\}$,
      then $\ell_2$ and~$\ell_3$ together contain 3~colors, and we are done by Case~2.
      So $\chi(\ell_3)$ is either $\{c\}$ or~$\{a,c\}$. In the former case we clearly
      have rainbow triangles.\footnote{{Of every  area multiple of $1/2$.}}
      So suppose $\chi(\ell_3)=\{a,c\}$.
      We shall show that if we have no rainbow triangle of area~$1/2$ with vertices in~$\ell_1\cup\ell_2\cup\ell_3$,
      then we can find monochromatic rectangles.\footnote{Of every area.}
      Fix consecutive vertical lines $\nu_1$, $\nu_2$ such that $\chi(\nu_1\cap\ell_1)=a$ and $\chi(\nu_2\cap\ell_1)=b$.
      If~$\chi(\nu_2\cap \ell_3)=c$, then $\nu_2$ contains 3~colors and we are done by Case~1 (as its argument
      works replacing horizontal lines by vertical lines); moreover, we can not have $\chi(\nu_1\cap \ell_3)=c$ else we
      obtain a rainbow triangle with vertices {$\nu_2\cap\ell_1$, $\nu_2\cap \ell_2$} and $\nu_1\cap \ell_3$.
      It follows that {$\chi(\nu_1\cap\ell_i)=a$} for each $i\in\{1,\,2,\,3\}$.

      To conclude, fix $d\in\N$ and a vertical line $\nu'$ at distance $d$ from $\nu_1$;
      we shall find a monochromatic rectangle of area~$d$ or a rainbow triangle.
      If $\chi(\nu'\!\cap\ell_1)=a$, then $\nu_1\cap \ell_1,\,\nu'\!\cap \ell_1,\,\nu_1\cap \ell_2$ and
      ~$\nu'\!\cap \ell_2$ are the vertices
      of a monochromatic rectangle of area~$d$. And if $\chi(\nu'\cap \ell_1)=b$, then as before we must have $\chi(\nu'\cap \ell_3)=a$
      (as otherwise the line $\nu'$ contains~3 distinct colors and we are done by Case~1).
      Hence $\nu_1\cap\ell_2,\,\nu'\!\cap\ell_2,\,\nu_1\cap\ell_3$, and~$\nu'\!\cap\ell_3$ form a monochromatic rectangle of area~$d$ as required.

      \emph{Case 4.} None of the conditions in the three above cases are satisfied. Then every three consecutive lines contain only at most two colors.  Let $L$ be a maximal set of consecutive lines containing only two colors $a$ and $b$, say.  Then we find three consecutive lines $\ell_1,\,\ell_2$, and~$\ell_3$, such that $\ell_1,\ell_2\in L$, $\ell_3\not\in L$, and $\ell_3$ contains a third color, say~$c$. Since no set of three consecutive lines out of these three contains $3$~colors, we have that $\ell_1$ and~$\ell_2$ are monochromatic of the same color. Hence we find monochromatic rectangles of every integer area.
    \end{proof}

\end{document}
